\begin{large}
// Source: https://github.com/malfple/quickpastetemplate/ \\
A special case of boolean satisfiability. Can a boolean equation equates to true?

$f = \left(x_1 \lor y_1 \right) \land \dots \land \left(x_n \lor y_n \right)$ ; $x_i \lor y_i \equiv \neg x_i \to y_i \equiv \neg y_i \to x_i$

Create a graph with 2n vertices, for each i, make 2 directed edges from $\neg x_i$ to $y_i$ and from $\neg y_i$ to $x_i$. $f$ is not satisfiable if both $x_i$ and $\neg x_i$ is in the same SCC (see Tarjan above). Note that for $\oplus$, you can convert it into the above equation. $x \oplus y = \left(x \lor y \right) \land \left( \neg x \lor \neg y \right)$ \\

\textbf{Backtracking} \\
if $x_i \to \neg x_i$ then $x_i$ has to be \textbf{FALSE}. This is valid for the opposite. This is also why if they are in the same SCC, a contradiction occurs. To assign values to $x_i$, you can use topological sort. Anything that comes earlier has to be FALSE. In the above implementation, you can use scc[i] (see NOTE). Sometimes you can just \texttt{if(scc[i] > scc[-i]) i = false, else i = true}

If in doubt: You need to record the members of each scc as an addition to the tarjan implementation below. \\
\begin{lstlisting}
iterate i = sccid..1:
if scc[i] doesn't have value, set to false
set the negation of the members of scc[i] to their respective value. The scc that they're in, now have a value.
if scc[i] has value true, all of its children has to be true
\end{lstlisting}
\end{large}
